\section{Introduction}

Recall from \textit{IA Numbers and Sets} that a set \(X\) is \term{countable} if \(X = \emptyset\), or there is some surjection \(\functiondomain{f}{\NN}{X}\).

We recall the following facts about countable sets:
\begin{itemize}
    \item Countable union of countable sets is countable. (We will refer to this as \term{zig-zagging}.)
    \item Diagonal argument will show sets that are uncountable.
    \item If \(X\) and \(Y\) are countable, then \(X \times Y\) are countable.
    \item If \(X\) is countable, then the set of \textbf{finite} subsets of \(X\) is countable.
    \item If \(X\) is infinite, then the \term{power set} of \(X\), denoted \(\powerset\pqty{X}\), is uncountable.
\end{itemize}

Now we introduce some definitions.

\begin{definition}[Alphabet]
    An \term{alphabet} \(\Sigma\) is a finite, non-empty set of symbols.
\end{definition}

We think of elements of \(\Sigma\) as \term{symbols} or \term{letters}.

\begin{examples*}
    \begin{itemize}
        \item The alphabet \(\Sigma = \Bqty{0, 1}\) is the binary alphabet.
        \item The alphabet \(\Sigma = \Bqty{0, 1, \ldots, 9}\) is the decimal alphabet.
        \item The alphabet \(\Sigma = \Bqty{a, b, c, \ldots, z}\) is the lowercase English alphabet.
        \item Suppose we have a group \(G\) and \(x, y \in g\). Then (ignoring the degenerate cases) we could take an alphabet of size \(4\) given by \(\Sigma = \Bqty{x, y, x\inverse, y\inverse} = \Bqty{x, y, X, Y}\) where \(X = x\inverse\), \(Y = y\inverse\) are \term{formal inverses} of \(x\) and \(y\) respectively.
    \end{itemize}
\end{examples*}

\begin{definition}[String/Word]
    A \term{string} or \term{word} of length \(m\) over an alphabet \(\Sigma\) is given by a list of length \(m\) (more formally, an \(m\)-tuple)
    \[
        \alpha = x_1 x_2 \ldots x_m
    \]
    where \(x_i \in \Sigma\). We denote the \term{length} of \(\alpha\) by \(\abs{\alpha} = m\). The \term{empty word} is denoted by \(\emptyword\) and has length \(0\). In other words,
    \[
        \Sigma^0 = \Bqty{\emptyword} \qand \abs{\emptyword} = 0.
    \]

    Let \(\Sigma\kstar\) denote the set of \term{all strings over \(\Sigma\)}, i.e.
    \[
        \Sigma\kstar = \bigcup_{n \in \NN} \Sigma^n,
    \]
    and let \(\Sigma\kplus\) denote the set of \term{all non-empty strings over \(\Sigma\)}, i.e.
    \[
        \Sigma\kplus = \Sigma\kstar \setminus \Bqty{\emptyword}.
    \]
\end{definition}

\begin{remark}
    Note the number of words of length \(n\) is give by
    \[
        \abs{\Sigma^n} = \abs{\Sigma}^n,
    \]
    and the set \(\Sigma\kstar\) is countable.
\end{remark}

\begin{definition}[Concatanation, Substring/Subword]
    If \(\alpha = x_1 x_2 \ldots x_m \in \Sigma^m\) and \(\beta = y_1 y_2 \ldots y_n \in \Sigma^n\), then the \term{concatenation} \(\alpha \beta\) is \(x_1 \ldots x_m y_1 \ldots y_n \in \Sigma^{m + n}\), with \(\abs{\alpha \beta} = \abs{\alpha}\abs{\beta}\).

    For \(0 \leq i \leq j \leq m\), \(x_{i + 1} \ldots x_j\) is a \term{substring/word} of \(\alpha\). In particular, if \(i = j\), then \(x_{i + 1} \ldots x_j = \emptyword\) is the empty string.
\end{definition}

\begin{remark}
    In other words, for \(\alpha, \beta \in \Sigma\kstar\), \(\beta\) is a substring of \(\alpha\) if there exist \(\gamma, \delta \in \Sigma\kstar\) such that \(\alpha = \gamma \beta \delta\).
\end{remark}

\begin{definition}[Homomorphism]
    For alphabets \(\Sigma\) and \(\Pi\), a \term{homomorphism} is a function \(\functiondomain{h}{\Sigma\kstar}{\Pi\kstar}\) such that for all \(\alpha, \beta \in \Sigma\kstar\),
    \[
        h\pqty{\alpha \beta} = h\pqty{\alpha} h\pqty{\beta}.
    \]

    In particular, \(h\pqty{\emptyword} = \emptyword\).
\end{definition}

\begin{remark}
    Alternatively, given any function \(\functiondomain{f}{\Sigma}{\Pi\kstar}\), there exists a unique homomorphism extending \(f\) by concatenation (by induction on the word length).
\end{remark}

\begin{definition}[Language]
    A \term{(formal) language} \(L\) over an alphabet \(\Sigma\) is a subset of \(\Sigma\kstar\).
\end{definition}

\begin{remark}
    Any language \(L\) is countable since we are picking it from a countable set. However, the set of all languages over \(\Sigma\) is uncountable since it is the power set of \(\Sigma\kstar\).
\end{remark}

\begin{examples*}
    \begin{itemize}
        \item Suppose \(\Sigma = \Bqty{a, b}\). Then a language over \(\Sigma\) could be given by
        \[
            L = \Bqty{\text{all words with the same number of \(a\)'s and \(b\)s}}.
        \]
        \item Suppose \(\Sigma = \Bqty{1}\), then we can think of a language over \(\Sigma\) as \(L \subseteq \NN = \Sigma\kstar\).
        \item Suppose \(\Sigma = \Bqty{a, b, \ldots, z}\). Then \(L\) could be the set of all words that are valid English words. In this case, \(L\) is a \term{finite language} since there are only finitely many valid English words.
        \item Let \(G = \angBrac{x_1, \ldots, x_n}\) be a finitely-generated group. Then with the alphabet
        \[
            \Sigma = \Bqty{x_1, \ldots, x_n, X_1, \ldots, X_n},
        \]
        then a word \(w \in \Sigma\kstar\) becomes an element of \(G\) by multiplying out. Define the language
        \[
            L = \set{w \in \Sigma\kstar}{w = e \text{ in } G}.
        \]

        Say \(G = \ZZ = \angBrac{x}\) is the infinite cyclic group (\(\ZZ\)). Then
        \[
            L = \Bqty{\text{all words with the same number of \(x\)'s and \(X\)s}}.
        \]
    \end{itemize}
\end{examples*}